\section{Introduction}
\label{sec:introduction}

This is the third presentation of Vibe Theory. The first (Pollard, 2025) set out the framework as philosophy. The second (Pollard, 2026) gave it a discrete mathematical structure built entirely from experience, with a concrete simulable target. This paper takes that target and runs it. It reports what a finite, reproducible simulation of the model actually shows, and it assembles those results into one picture of how space, time, matter, force, gravity, and quantum behaviour fit together in a universe made of one substance.

The central claim of Vibe Theory is monism. There is one kind of thing, the vibe, a state that stands in relation to other states. There is no background space and no background time. The only primitive relation is the note. Everything else, geometry and matter and force, is a large-scale pattern of the vibe mesh. A monist theory of this strictness owes a debt that most ontologies never pay. It must show that the variety of the physical world can actually come out of one substrate, not merely assert that it could. This paper is an attempt to pay that debt in simulation.

The method is deliberately empirical about the model, even though the model itself is not yet empirically tested against nature. We treat each foundational question as a measurement. We built a testbed, a finite and seeded program with a known-answer test suite, and posed nine questions that the framework must answer to be coherent: whether the dynamics can have a stable ground state and a local time, whether a smooth spacetime is dynamically stabilized rather than assumed, whether a substrate can be at once Lorentz-safe and navigable, whether matter and spin are patterns of the mesh, whether the recovered geometry is unique, whether the sum over histories is computable and lands on the right dimension, whether quantum statistics can come from a deterministic base, whether gauge force and a charged fermion can be built, and where the framework's own boundary lies. Each question becomes a runnable experiment, and each experiment returns a number.

The results are mixed in the honest way a real research program is mixed, and we present them that way. Several questions are answered cleanly and validated against known-answer cases. Two of the hardest, the dynamical stabilization of a spacetime phase and the resolution of the law, are addressed at the scale we can reach, the first showing a stable phase and a first-order signature with strict dominance still open, the second by relocating the Hamiltonian from the microscopic rule to the emergent mesh. One, the quantum link, is made precise rather than closed. A few are genuinely open in physics itself, or lie at the framework's philosophical boundary, and we mark them as such rather than overclaim.

Throughout, the framework sits in the family of discrete-physics and digital-physics programs. The substrate is a causal set in the sense of Bombelli, Lee, Meyer, and Sorkin (1987), the dynamics borrows the Benincasa-Dowker action (2010), the matter sector uses lattice fermions and the overlap construction of Neuberger (1998) and Ginsparg and Wilson (1982), the gauge sector uses Wilson's lattice gauge theory (1974), and the quantum question engages 't Hooft's Cellular Automaton Interpretation (2016) and Bell's theorem (1964). What Vibe Theory adds is the insistence that all of these live on one experiential substrate, and the demonstration, here, that they can.

This work is a philosophical and mathematical model with a simulable build target rather than a finished scientific theory. References to physics are sources of structure and tests of coherence, not claims to have derived nature. What the paper offers is a candidate ontology that has been put under simulation, with its successes, its partial results, and its open frontiers all stated precisely.
